\documentclass[11pt]{article}
\usepackage[preprint]{acl}
\usepackage{natbib}
\usepackage{times}
\usepackage{latexsym}
\usepackage[T1]{fontenc}
\usepackage[utf8]{inputenc}
\usepackage{microtype}
\usepackage{graphicx}
\usepackage{booktabs}

\title{An Empirical Analysis of Hallucination Behavior in LLMs in the Medical Domain}

\author{Aidan Perez \\
  MS in AI, Boston University\\
  \texttt{aperez00@bu.edu} \\\And
  Dang Dinh \\
  MS in DS, Boston University\\
  \texttt{dddinh@bu.edu} \\\AND
  Kaveri Goel \\
  MS in AI, Boston University\\
  \texttt{kvrgl@bu.edu} \\\And
  Ying Huang \\
  MS in CS, Boston University\\
  \texttt{yinghy@bu.edu} }

\begin{document}
\maketitle

% ── ABSTRACT ─────────────────────────────────────────────────────────────────
\begin{abstract}
Large Language Models (LLMs) have demonstrated remarkable capabilities 
in generating fluent, contextually relevant text. However, they 
frequently produce hallucinations --- outputs that appear plausible but 
are factually incorrect, unsupported, or fabricated. This problem is 
especially critical in high-stakes domains such as medicine, where 
incorrect information can have serious real-world consequences. We 
present a controlled empirical study of hallucination behavior across 
three model families (Llama-3-8B, BioMistral-7B, and BioGPT) on two 
medical QA benchmarks (MedQA and PubMedQA), evaluating four 
intervention conditions: zero-shot prompting, chain-of-thought (CoT) 
prompting, retrieval-augmented generation (RAG), and supervised 
fine-tuning (SFT). Our results show that larger instruction-tuned 
models benefit most from RAG, CoT prompting consistently degrades 
performance relative to zero-shot due to output noise and parse 
failures, and fine-tuning improves BioGPT accuracy by 6.66 percentage 
points while eliminating parse failures entirely. Combining fine-tuning 
with RAG, however, leads to catastrophic output degradation, suggesting 
that task-specific fine-tuning and retrieval augmentation are 
incompatible in their current forms for this model.
\end{abstract}

% ── INTRODUCTION ─────────────────────────────────────────────────────────────
\section{Introduction}

Large Language Models (LLMs) have revolutionized natural language 
processing by demonstrating unprecedented capabilities in generating 
coherent, contextually relevant text. However, a significant challenge 
is their propensity to produce hallucinations --- outputs that may be 
fluent and plausible but are factually incorrect, unsupported by 
evidence, or entirely fabricated. This issue is particularly concerning 
in medicine, where incorrect information can lead to serious consequences 
for patient care and safety.

Despite the growing body of research on hallucination in general NLP 
contexts, relatively little empirical work systematically characterizes 
hallucination behavior in medical LLMs, or directly compares the 
effectiveness of inference-time versus training-time mitigation 
strategies. In this paper, we address this gap by evaluating multiple 
LLMs on standardized medical QA benchmarks under four experimental 
conditions: zero-shot prompting, chain-of-thought prompting, 
retrieval-augmented generation, and supervised fine-tuning. By 
systematically comparing these interventions, we aim to answer not only 
how often LLMs hallucinate on medical questions, but which type of 
intervention works best --- and whether combining them yields further 
gains.

% ── RELATED WORK ─────────────────────────────────────────────────────────────
\section{Related Work}

\subsection{Hallucination in Large Language Models}
Hallucination is a widely recognized problem in LLMs, where a model 
generates fluent and plausible text that is factually incorrect, 
unsupported by the input, or inconsistent with external knowledge. Prior 
work has categorized hallucinations into intrinsic hallucinations, which 
contradict the source context, and extrinsic hallucinations, which 
introduce information not supported by the provided evidence 
\cite{ji2023survey}. Hallucination has also been studied in abstractive 
summarization, where models often generate coherent but unfaithful 
content \cite{maynez2020faithfulness}. These findings are relevant to 
medical LLMs because medical responses require both factual accuracy and 
faithfulness to clinical context.

\subsection{Medical Question Answering and Clinical Reasoning}
Recent work has shown that LLMs can achieve strong performance on 
medical QA benchmarks such as MedQA, MedMCQA, and PubMedQA, but remain 
vulnerable to confident errors on complex clinical reasoning tasks 
\cite{lievin2024medical}. Medical benchmarks provide a useful setting 
for studying hallucination because they offer verifiable ground-truth 
answers. Prior studies have also developed medical hallucination 
benchmarks, such as Med-HALT, which evaluates reasoning-based and 
memory-based hallucination in medical LLMs \cite{pal2023medhalt}.

\subsection{Chain-of-Thought Prompting}
Chain-of-thought (CoT) prompting encourages a model to generate 
intermediate reasoning steps before producing a final answer 
\cite{wei2022chain}. While CoT has been shown to improve reasoning 
on complex tasks, it can also produce longer generations that introduce 
additional opportunities for error. In medical QA, a model may generate 
a fluent but incorrect reasoning chain before giving its final answer, 
and CoT outputs may be harder to parse reliably \cite{wang2023selfconsistency}.

\subsection{Retrieval-Augmented Generation}
Retrieval-augmented generation (RAG) combines a generative model with 
an external retrieval component, allowing the model to condition its 
output on retrieved documents rather than relying only on parametric 
knowledge \cite{lewis2020retrieval}. In the medical domain, retrieval 
can help ground model outputs in biomedical abstracts or clinical 
documents, potentially reducing extrinsic hallucination. However, RAG 
does not fully eliminate hallucination, as a model may still 
misinterpret retrieved evidence or fail to align its answer with the 
retrieved context.

\subsection{Fine-Tuning for Medical NLP}
Supervised fine-tuning (SFT) on domain-specific data has been shown to 
improve model performance on medical tasks \cite{singhal2023large}. 
Domain-specific pretraining, as in BioGPT \cite{luo2022biogpt} and 
BioMistral \cite{labrak2024biomistral}, has also been proposed as a way 
to reduce hallucination by internalizing medical knowledge during 
training. However, the relative effectiveness of fine-tuning versus 
inference-time interventions such as RAG remains underexplored, 
motivating our direct comparison.

% ── METHODS ──────────────────────────────────────────────────────────────────
\section{Methods}

\subsection{Datasets}
We evaluate on two medical QA benchmarks that represent different task 
types and hallucination detection modalities.

\textbf{MedQA (USMLE)} \cite{jin2021disease} consists of approximately 
12,000 multiple-choice questions derived from United States Medical 
Licensing Examinations, covering a broad range of clinical topics. 
Gold-standard answer letters allow direct accuracy measurement. We 
evaluate on the held-out test split, sampling 450--500 items per 
experimental condition.

\textbf{PubMedQA} \cite{jin2019pubmedqa} consists of 450 
expert-annotated yes/no/maybe questions derived from PubMed research 
abstracts. Each question is paired with the source abstract, making it 
ideal for evaluating context-grounded generation and retrieval 
augmentation. We use the full labeled training split (450 items) for 
all PubMedQA conditions.

\subsection{Models}
We compare three models spanning capability tiers:

\textbf{Meta-Llama-3-8B-Instruct} is a large general-purpose 
instruction-tuned model serving as our strong baseline. It was not 
fine-tuned on medical data.

\textbf{BioMistral-7B} \cite{labrak2024biomistral} is a 
medically-adapted model built on Mistral-7B and further pretrained on 
biomedical corpora. It allows us to test whether domain-specific 
pretraining reduces hallucination relative to a general model.

\textbf{BioGPT} \cite{luo2022biogpt} is a smaller (1.5B parameter) 
generative model pretrained on PubMed abstracts. It serves as both a 
small-model baseline and the target for our supervised fine-tuning 
experiments.

\subsection{Experimental Conditions}
We evaluate each model under the following conditions:

\textbf{Zero-shot prompting} provides the model with the question and 
answer choices (MedQA) or question alone (PubMedQA) with no additional 
context or reasoning instruction.

\textbf{Chain-of-thought (CoT) prompting} asks the model to reason 
step-by-step before providing a final answer, ending with 
\texttt{Answer: <letter>}.

\textbf{Retrieval-augmented generation (RAG)} uses BM25 retrieval over 
the PubMedQA abstract corpus with top-$k{=}2$ abstracts prepended as 
context, using a fixed prompt template across all models.

\textbf{Supervised fine-tuning (SFT)} fine-tunes BioGPT on the full 
MedQA training split (10,178 samples) for 3 epochs using standard 
next-token prediction loss, learning rate $2 \times 10^{-5}$, effective 
batch size 16, and fp16 mixed precision on a single A100 GPU.

\textbf{SFT + RAG} applies RAG at inference time to the fine-tuned 
BioGPT model, testing whether the two interventions are complementary 
or redundant.

\subsection{Evaluation Metrics}

\textbf{Accuracy} is computed as the proportion of outputs whose 
extracted answer exactly matches the gold-standard answer. For MedQA, 
this means the extracted letter (A--D) matches the gold answer letter. 
For PubMedQA, this means the extracted decision (yes/no/maybe) matches 
the gold label.

\textbf{Error rate} is defined as $1 - \text{accuracy}$ over all 
outputs including parse failures. For MedQA conditions (zero-shot, CoT, 
SFT), we use error rate rather than ``hallucination rate'' since there 
is no retrieved context to assess faithfulness against --- incorrect 
answers reflect factual errors rather than hallucination in the 
retrieval faithfulness sense.

\textbf{Hallucination rate} is reserved for RAG conditions on PubMedQA, 
where it captures outputs that are either incorrect or unfaithful to the 
retrieved context. Parse failures are included in the hallucination 
count since an unparseable output is clinically unusable.

\textbf{Parse failure rate} is reported separately as the proportion of 
outputs from which no structured answer could be extracted using regex 
pattern matching. For MedQA, the extractor looks for \texttt{Answer: 
[A-D]} or a leading letter. For PubMedQA, it looks for yes, no, or 
maybe at the start of or within the output. Parse failures are counted 
as errors in all conditions.

\textbf{BERTScore F1} \cite{zhang2019bertscore} measures token-level 
semantic similarity between generated outputs and reference answers 
using contextual embeddings from RoBERTa-large. It serves as a 
secondary signal of output quality independent of exact-match accuracy.

% ── RESULTS ──────────────────────────────────────────────────────────────────
\section{Results}

\subsection{RQ1: Zero-Shot and Chain-of-Thought on MedQA}

Table~\ref{tab:medqa-llama} summarises the performance of 
Meta-Llama-3-8B-Instruct on MedQA under zero-shot and CoT prompting.

\begin{table}[h]
  \centering
  \begin{tabular}{lcc}
    \toprule
    \textbf{Metric} & \textbf{Zero-Shot} & \textbf{CoT} \\
    \midrule
    Correct           & 302       & 282      \\
    Accuracy          & 60.4\%    & 56.4\%   \\
    Error Rate        & 39.6\%    & 43.6\%   \\
    Parse Failures    & 4         & 18       \\
    BERTScore F1      & 0.819     & 0.814    \\
    \bottomrule
  \end{tabular}
  \caption{MedQA results for Llama-3-8B-Instruct (500 samples).}
  \label{tab:medqa-llama}
\end{table}

Zero-shot prompting outperforms CoT on every metric. CoT produces 
longer reasoning traces that introduce more formatting artifacts, 
repetition, and parse failures, leading to a higher measured error rate. 
These results suggest that explicit reasoning prompts do not 
automatically improve performance on medical multiple-choice QA.

We observe a consistent pattern with Qwen2.5-3B-Instruct on MedQA 
(200 samples), as shown in Table~\ref{tab:medqa-qwen}.

\begin{table}[h]
  \centering
  \begin{tabular}{lcc}
    \toprule
    \textbf{Metric} & \textbf{Zero-Shot} & \textbf{CoT} \\
    \midrule
    Correct           & 86        & 47       \\
    Accuracy          & 43.0\%    & 23.5\%   \\
    Error Rate        & 57.0\%    & 76.5\%   \\
    Parse Failures    & 1         & 12       \\
    \bottomrule
  \end{tabular}
  \caption{MedQA results for Qwen2.5-3B-Instruct (200 samples).}
  \label{tab:medqa-qwen}
\end{table}

Qwen2.5-3B-Instruct performs substantially worse than 
Llama-3-8B-Instruct under both conditions. The gap between zero-shot 
and CoT is even larger for Qwen2.5-3B, suggesting that smaller models 
are more susceptible to the noise introduced by chain-of-thought 
reasoning.

\subsection{RQ2: Retrieval-Augmented Generation on PubMedQA}

To directly evaluate whether RAG reduces hallucination, we ran both 
zero-shot and RAG conditions on the same 450 PubMedQA questions for all 
three models. Tables~\ref{tab:pubmedqa-zeroshot} and~\ref{tab:pubmedqa-rag} 
report the results.

\begin{table}[h!]
  \centering
  \small
  \begin{tabular}{lcccc}
    \toprule
    \textbf{Model} & \textbf{Acc.} & \textbf{Halluc.} & \textbf{Parse Fail} & \textbf{BERT} \\
    \midrule
    Llama-3-8B    & 56.00\% & 44.00\% &  0.00\% & 0.830 \\
    BioMistral-7B & 46.00\% & 54.00\% & 16.22\% & 0.838 \\
    BioGPT        & 22.89\% & 77.11\% & 57.56\% & 0.785 \\
    \bottomrule
  \end{tabular}
  \caption{PubMedQA zero-shot results (450 samples per model).}
  \label{tab:pubmedqa-zeroshot}
\end{table}

\begin{table}[h!]
  \centering
  \small
  \begin{tabular}{lcccc}
    \toprule
    \textbf{Model} & \textbf{Acc.} & \textbf{Halluc.} & \textbf{Parse Fail} & \textbf{BERT} \\
    \midrule
    Llama-3-8B    & 71.11\% & 28.89\% &  3.33\% & 0.853 \\
    BioMistral-7B & 45.11\% & 54.89\% & 14.67\% & 0.820 \\
    BioGPT        & 10.89\% & 89.11\% & 79.78\% & 0.801 \\
    \bottomrule
  \end{tabular}
  \caption{PubMedQA RAG results with BM25 top-2 retrieval (450 samples per model).}
  \label{tab:pubmedqa-rag}
\end{table}

The results reveal three distinct patterns across model tiers. For 
Llama-3-8B-Instruct, RAG substantially improves performance --- accuracy 
increases from 56.00\% to 71.11\% (+15.11 points) and hallucination 
rate drops from 44.00\% to 28.89\%, with BERTScore F1 also improving 
from 0.830 to 0.853. For BioMistral-7B, RAG has essentially no effect 
--- accuracy changes by only $-$0.89 points (46.00\% $\rightarrow$ 
45.11\%), suggesting that domain-pretrained models already possess the 
parametric knowledge that retrieval provides and do not benefit from 
additional context. For BioGPT, RAG actively degrades performance --- 
accuracy drops from 22.89\% to 10.89\% ($-$12 points) and parse 
failures increase from 57.56\% to 79.78\%, indicating that the retrieved 
context confuses the model rather than helping it generate structured 
responses.

\subsection{RQ3: Effect of Fine-Tuning on Hallucination}

We fine-tuned BioGPT on the full MedQA training split (10,178 samples) 
for 3 epochs. Training loss decreased steadily across epochs 
($2.05 \rightarrow 1.73 \rightarrow 1.63$) with validation loss 
following the same trend ($1.78 \rightarrow 1.72 \rightarrow 1.71$), 
confirming convergence without overfitting. The fine-tuned model was 
evaluated on a matched subsample of 450 MedQA test items (seed=42) 
alongside the zero-shot baseline. Table~\ref{tab:rq3} reports the 
results.

\begin{table}[h!]
  \centering
  \small
  \begin{tabular}{lccc}
    \toprule
    \textbf{Condition} & \textbf{Accuracy} & \textbf{Error Rate} & \textbf{Parse Fail} \\
    \midrule
    BioGPT zero-shot   & 17.78\% & 82.22\% & 23.11\% \\
    BioGPT fine-tuned  & 24.44\% & 75.56\% &  0.00\% \\
    \bottomrule
  \end{tabular}
  \caption{RQ3: BioGPT on MedQA before and after fine-tuning 
  (450 matched samples).}
  \label{tab:rq3}
\end{table}

Fine-tuning improved accuracy by 6.66 percentage points and 
completely eliminated parse failures (23.11\% $\rightarrow$ 0.00\%), 
indicating the model learned to produce clean answer letters 
consistently. Despite this improvement, fine-tuned BioGPT (24.44\%) 
remains substantially below zero-shot Llama-3-8B-Instruct (60.4\%), 
suggesting that model scale and instruction-tuning capability are more 
important factors than domain-specific fine-tuning alone.

\subsection{RQ4: RAG and Fine-Tuning Interaction}
 
We applied RAG at inference time to the fine-tuned BioGPT model on 
PubMedQA to test whether the two interventions are complementary. 
Table~\ref{tab:rq4} reports the results for three BioGPT conditions: 
RAG only (base model, PubMedQA), SFT only (fine-tuned, MedQA), and 
SFT + RAG (fine-tuned, PubMedQA with retrieval).
 
\begin{table}[h!]
  \centering
  \small
  \begin{tabular}{lccc}
    \toprule
    \textbf{Condition} & \textbf{Acc.} & \textbf{Halluc.} & \textbf{Parse Fail} \\
    \midrule
    BioGPT RAG only  & 10.89\% & 89.11\% & 79.78\% \\
    BioGPT SFT only  & 24.44\% & 75.56\% &  0.0\%  \\
    BioGPT SFT + RAG &  3.33\% & 96.67\% & 93.33\% \\
    \bottomrule
  \end{tabular}
  \caption{RQ4: BioGPT under RAG, SFT, and combined conditions. 
  RAG and SFT+RAG evaluated on PubMedQA (450 samples); 
  SFT evaluated on MedQA (450 samples).}
  \label{tab:rq4}
\end{table}
 
BioGPT struggles under all conditions on PubMedQA --- zero-shot 
achieves only 22.89\% (Table~\ref{tab:pubmedqa-zeroshot}), RAG alone 
drops this further to 10.89\%, and SFT + RAG collapses to 3.33\% with 
a 93.33\% parse failure rate. The consistent degradation across 
conditions suggests that BioGPT fundamentally lacks the capacity for 
structured yes/no/maybe output, independent of whether retrieval or 
fine-tuning is applied.

An important methodological caveat must be noted: the SFT model was 
fine-tuned on MedQA (multiple-choice, A/B/C/D format) and then 
evaluated on PubMedQA (yes/no/maybe format). These are fundamentally 
different output formats. The observed degradation in the SFT + RAG 
condition therefore reflects not only BioGPT's limited capacity, but 
also a \textit{cross-task format mismatch} --- the fine-tuned model 
learned to produce answer letters and lost the output flexibility 
required for yes/no/maybe responses. This confounding effect limits 
the conclusions that can be drawn about the RAG + fine-tuning 
interaction specifically.
 
Additionally, PubMedQA contains only 450 labeled samples in total, 
all of which were used for evaluation. This small dataset size limits 
statistical power and generalizability.

Despite these limitations, the result reveals a practically important 
finding: task-specific fine-tuning on a fixed output format introduces 
format rigidity that compounds existing model limitations, making 
the combined SFT + RAG approach worse than either intervention alone. 
A rigorous evaluation of the RAG + fine-tuning interaction would 
require either (1) fine-tuning on PubMedQA specifically and evaluating 
with and without RAG on a held-out split, or (2) using a larger dataset 
that supports both training and evaluation splits. We leave this for 
future work.

\subsection{Prediction Distribution Analysis}

Qwen2.5-3B-Instruct shows a strong bias toward hedging. On PubMedQA, 
87.5\% of its predictions were ``maybe,'' compared to a gold-label 
``maybe'' rate of only 11.0\%. On MedQA (zero-shot), the model 
exhibited a preference for option~D (55.0\% of predictions vs.\ 26.0\% 
in the gold labels), while CoT prompting shifted the distribution 
heavily toward option~A (60.5\%). This sensitivity of the answer 
distribution to the prompting strategy suggests that the model's 
predictions are driven more by surface-level prompt cues than by 
genuine medical reasoning.

% ── DISCUSSION ───────────────────────────────────────────────────────────────
\section{Discussion}
 
\subsection{Effect of Model Scale}
Across all conditions, model scale and instruction-tuning capability 
emerge as the dominant factors in hallucination reduction. 
Llama-3-8B-Instruct consistently outperforms both BioMistral-7B and 
BioGPT despite having no medical domain pretraining. This is consistent 
with recent findings that general large models with strong instruction 
tuning can match or exceed smaller specialized models on medical 
benchmarks \cite{singhal2023large}. Our results suggest that investing 
in larger, better-aligned general models may be more effective than 
domain-specific pretraining for reducing hallucination in medical QA.
 
\subsection{Inference-Time vs.\ Training-Time Interventions}
Our results reveal an important asymmetry between inference-time and 
training-time interventions. RAG substantially reduces hallucination for 
Llama-3-8B (+15.11 accuracy points, hallucination drops from 44.00\% to 
28.89\%) but has no measurable effect on BioMistral-7B ($-$0.89 points) 
and actively hurts BioGPT ($-$12 points). This suggests that RAG is 
only effective when the model has sufficient instruction-following 
capacity to leverage retrieved context. Fine-tuning improves BioGPT 
accuracy by 6.66 percentage points and completely eliminates parse 
failures (23.11\% $\rightarrow$ 0.00\%), demonstrating that 
task-specific adaptation is beneficial even for small models. However, 
the absolute accuracy of the fine-tuned model (24.44\%) remains far 
below zero-shot Llama-3-8B (60.4\%), suggesting that model scale and 
instruction-tuning capability are more critical factors than fine-tuning 
alone for hallucination reduction in medical QA.
 
\subsection{Format Rigidity as a Hidden Risk of Fine-Tuning}
A particularly important finding emerges from our RQ4 experiment. 
BioGPT struggles on PubMedQA under all conditions --- zero-shot 
(22.89\%), RAG (10.89\%), and SFT + RAG (3.33\%) --- revealing that 
the model fundamentally lacks the capacity for structured yes/no/maybe 
output. Critically, applying SFT compounds this limitation: the model, 
having been trained exclusively on MedQA multiple-choice format, 
further loses output flexibility and achieves a 93.33\% parse failure 
rate under SFT + RAG.

This finding has practical implications beyond our specific experimental 
setup: fine-tuning a model on a fixed output format introduces format 
rigidity that may prevent it from generalizing to other tasks at 
inference time, and can compound existing model weaknesses rather than 
mitigate them. Practitioners deploying fine-tuned models in multi-task 
clinical settings should carefully evaluate format compatibility before 
combining fine-tuning with retrieval augmentation.

We acknowledge an important caveat: our RQ4 setup involves a 
cross-task format mismatch, and PubMedQA's 450 labeled samples are 
too few to support both a fine-tuning split and a reliable evaluation 
split. A fully controlled interaction test --- fine-tuning and RAG 
evaluation within the same format and dataset --- remains for future 
work. Nevertheless, the finding that SFT compounds rather than 
alleviates BioGPT's structured output failures stands independently 
of this limitation.

% ── CONCLUSION ───────────────────────────────────────────────────────────────
\section{Conclusion}
 
We presented a controlled empirical study of hallucination behavior in 
medical LLMs across four intervention conditions. Our key findings are: 
(1) larger instruction-tuned models hallucinate less under RAG, with 
Llama-3-8B achieving 71.11\% accuracy on PubMedQA; (2) CoT prompting 
consistently degrades performance relative to zero-shot due to output 
noise and parse failures; (3) fine-tuning BioGPT improves accuracy by 
6.66 percentage points and completely eliminates parse failures, though 
absolute accuracy remains bounded by model scale; and (4) combining 
fine-tuning with RAG reveals a format rigidity problem --- task-specific 
fine-tuning can severely degrade output format flexibility, preventing 
generalization to new task formats even with retrieved context. These 
findings highlight that model scale and instruction-tuning capability 
are more important than domain-specific adaptation for hallucination 
reduction, and that combined interventions must be carefully evaluated 
for format compatibility rather than assumed to be additive.

% ── LIMITATIONS ──────────────────────────────────────────────────────────────
\section{Limitations}
 
Several limitations of this work should be noted. First, our RQ4 
evaluation involves a cross-task format mismatch: BioGPT was fine-tuned 
on MedQA (multiple-choice, A/B/C/D) but evaluated with RAG on PubMedQA 
(yes/no/maybe). A rigorous interaction test would require fine-tuning 
and evaluation within the same format and dataset. PubMedQA's 450 
labeled samples are too few to support both a fine-tuning split and a 
reliable evaluation split simultaneously, which is why we could not 
conduct this experiment within the current study.
 
Second, our fine-tuning used only 3 epochs with standard SFT; longer 
training or alternative objectives such as RLHF or DPO may yield 
greater hallucination reductions. Finally, we use BM25 retrieval with 
top-2 abstracts; denser retrieval methods such as MedCPT embeddings 
may improve RAG quality and faithfulness.

% ── REFERENCES ───────────────────────────────────────────────────────────────
\bibliography{references}

\end{document}